% ===== PRÉAMBULE CANONIQUE — à coller VERBATIM en tête de CHAQUE .tex =====
\documentclass[11pt,a4paper]{article}
\usepackage[utf8]{inputenc}\usepackage[T1]{fontenc}\usepackage{lmodern}
\usepackage[french]{babel}
\usepackage{amsmath,amssymb,mathtools}\usepackage{icomma}
\usepackage[version=4]{mhchem}          % \ce{...} : équations & formules
\usepackage{siunitx}\sisetup{locale=FR} % \qty{}{}, \si{}, \ang{}
\usepackage{chemfig}                    % \chemfig{...} : structures développées
\usepackage{xcolor}\usepackage{colortbl}
\usepackage{tikz}\usepackage{pgfplots}\pgfplotsset{compat=1.18}
\usepackage[most]{tcolorbox}\usepackage{enumitem}\usepackage{array,booktabs}
\usepackage[a4paper,margin=2.2cm]{geometry}
\usetikzlibrary{arrows.meta,calc,angles,quotes,positioning,patterns,backgrounds,shapes.geometric}
\definecolor{primary}{RGB}{222,49,99}\definecolor{primarydark}{RGB}{150,22,55}
\definecolor{defcol}{RGB}{0,114,178}\definecolor{methcol}{RGB}{52,92,148}
\definecolor{excol}{RGB}{0,135,90}\definecolor{attncol}{RGB}{200,40,55}
\tcbset{boxbase/.style={enhanced,breakable,boxrule=0.9pt,arc=2.5pt,
  left=3mm,right=3mm,top=2.3mm,bottom=2.3mm,fonttitle=\bfseries,coltitle=white}}
% --- boîtes TUTORIEL (noms EXACTS, ne pas renommer) ---
\newtcolorbox{cle}[1][]{boxbase,colback=defcol!6,colframe=defcol,
  title={\textsc{À retenir}\ifx\relax#1\relax\relax\else\textmd{ : \itshape #1}\fi}}
\newtcolorbox{methode}[1][]{boxbase,colback=methcol!7,colframe=methcol,sharp corners,
  title={\textsc{Méthode}\ifx\relax#1\relax\relax\else\textmd{ --- \itshape #1}\fi}}
\newtcolorbox{exemplebox}[1][]{boxbase,colback=excol!5,colframe=excol,
  title={\textsc{Exemple}\ifx\relax#1\relax\relax\else\textmd{ : \itshape #1}\fi}}
\newtcolorbox{piege}[1][]{boxbase,colback=attncol!6,colframe=attncol,
  title={\textsc{Piège}\ifx\relax#1\relax\relax\else\textmd{ : \itshape #1}\fi}}
\newtcolorbox{remarque}[1][]{boxbase,colback=black!4,colframe=black!45,
  title={\textsc{Remarque}\ifx\relax#1\relax\relax\else\textmd{ : \itshape #1}\fi}}
% --- boîtes EXERCICE / CORRIGÉ (noms EXACTS, auto-comptées) ---
\newtcolorbox[auto counter]{exo}[1][]{boxbase,colback=primary!4,colframe=primary,
  title={\textsc{Exercice}~\thetcbcounter\ifx\relax#1\relax\relax\else\textmd{ --- \itshape #1}\fi}}
\newtcolorbox[auto counter]{corrige}[1][]{boxbase,colback=excol!4,colframe=excol,
  title={\textsc{Corrigé}~\thetcbcounter\ifx\relax#1\relax\relax\else\textmd{ --- \itshape #1}\fi}}
\newcommand{\R}{\mathbb{R}}\newcommand{\N}{\mathbb{N}}\newcommand{\Z}{\mathbb{Z}}\newcommand{\ds}{\displaystyle}
% (un \usepackage{fancyhdr}+en-tête est possible ; il est ignoré côté web)
% =========================================================================
\begin{document}
\sisetup{per-mode=symbol}

\begin{exo}[repérer la (les) paire(s) qui n'est (ne sont) pas un couple conjugué]
Un couple acide/base conjugués relie deux espèces qui diffèrent d'\textbf{exactement un} \ce{H+}. Parmi les
paires suivantes, indique celle(s) qui \textbf{ne forme(nt) pas} un couple conjugué, et justifie.
\begin{enumerate}[label=\alph*)]
  \item \ce{HClO / ClO-}
  \item \ce{H3O+ / OH-}
  \item \ce{H2PO4- / HPO4^{2-}}
  \item \ce{NH4+ / NH2-}
  \item \ce{HSO4- / SO4^{2-}}
\end{enumerate}
\end{exo}

\begin{exo}[quelles réactions sont des réactions acide-base ?]
Parmi les réactions suivantes en milieu aqueux, indique lesquelles sont des \textbf{réactions acide-base}
(transfert de proton \ce{H+}) et non des précipitations. Justifie brièvement.
\begin{enumerate}[label=\arabic*)]
  \item \ce{HCN + NaOH -> NaCN + H2O}
  \item \ce{AgNO3 + NaCl -> AgCl v + NaNO3}
  \item \ce{NH4Cl + NaOH -> NH3 + NaCl + H2O}
  \item \ce{CuSO4 + 2 NaOH -> Cu(OH)2 v + Na2SO4}
  \item \ce{NaHCO3 + HCl -> NaCl + H2O + CO2 ^}
\end{enumerate}
\end{exo}

\begin{exo}[électrolyte fort, faible ou nul ?]
On dissout séparément dans l'eau : le glucose \ce{C6H12O6}, l'éthanol \ce{C2H5OH}, le chlorure de potassium
\ce{KCl}, et l'acide fluorhydrique \ce{HF}.
\begin{enumerate}[label=\alph*)]
  \item Lequel est un \textbf{électrolyte fort} (totalement dissocié en ions) ?
  \item Laquelle des quatre solutions, préparées à même concentration, conduit le \textbf{moins} le courant ?
        Pourquoi ?
\end{enumerate}
\end{exo}

\begin{exo}[mesurer un degré d'ionisation, puis le prévoir]
On prépare \qty{100}{mL} de solution d'acide acétique \ce{CH3COOH} à \qty{0,10}{mol\per L}. À l'équilibre, un
dosage révèle $\qty{5,0e-4}{mol}$ d'ions acétate \ce{CH3COO-}.
\begin{enumerate}[label=\alph*)]
  \item Calcule le degré d'ionisation $\alpha$.
  \item Cette solution est-elle un électrolyte fort ou faible ?
  \item Si l'on diluait cette solution \textbf{4 fois}, dans quel sens et de quel facteur (approché) varierait
        $\alpha$ ? (Loi de dilution d'Ostwald $\alpha \approx \sqrt{K_a/C_a}$.)
\end{enumerate}
\end{exo}

\begin{exo}[vrai ou faux sur les définitions]
Pour chaque affirmation, réponds par \textbf{vrai} ou \textbf{faux} et \textbf{justifie} en une ligne.
\begin{enumerate}[label=\alph*)]
  \item L'acide conjugué de l'ion \ce{HS-} est \ce{H2S}.
  \item La base conjuguée de l'eau \ce{H2O} est l'ion \ce{H3O+}.
  \item L'ion \ce{SO4^{2-}} est l'acide conjugué de \ce{HSO4-}.
  \item Tout acide de Lewis est nécessairement un acide de Brønsted.
  \item \ce{H3O+} ne possède pas d'acide conjugué observable en solution aqueuse.
\end{enumerate}
\end{exo}

\end{document}
